\documentclass[../main.tex]{subfiles}
\begin{document}

\section{Upload and Download traffic sources}
\label{app:upload-download-traffic}

Chapter~\ref{sec:down_up_program} describes Upload and Download at the S7comm protocol level. This appendix records why live block-transfer traffic was not captured in the lab and which reference PCAP source was used for rule development instead.

Attempts to simulate program upload and download in the available tooling were unsuccessful because of limitations in the simulation software.

When uploading program code from OpenPLC Editor to OpenPLC Runtime, the program is sent through the OpenPLC Runtime RESTful API rather than using the standard S7comm structure.

Figure~\ref{fig:ksdownupprogramimage5png} shows uploading a program from OpenPLC Editor to OpenPLC Runtime via RESTful API.

\begin{figure}[htbp]
    \centering
    \includegraphics[width=0.9\textwidth]{../attacks/down_up_program/image-5.png}
    \caption{Uploading a program from OpenPLC Editor to OpenPLC Runtime via RESTful API}
    \label{fig:ksdownupprogramimage5png}
\end{figure}

OpenPLC Runtime does not support loading programs in the usual Siemens S7comm way.

Another approach is to simulate program transfer on a Siemens PLC through PLCSIM Advanced. In that setup, the program would use the standard S7comm block-transfer sequence.

\textbf{Step 1} --- Set up a virtual switch for Siemens virtual devices:

Figure~\ref{fig:ksdownupprogramimage2png} shows virtual switch setup for Siemens virtual devices.

\begin{figure}[htbp]
    \centering
    \includegraphics[width=0.9\textwidth]{../attacks/down_up_program/image-2.png}
    \caption{Virtual switch setup for Siemens virtual devices}
    \label{fig:ksdownupprogramimage2png}
\end{figure}

\textbf{Step 2} --- Run a Siemens PLC in PLCSIM Advanced on the same network as the virtual switch. PLCSIM Advanced only supports simulating Siemens S7-1500 PLCs.

Figure~\ref{fig:cksdownupprogramimagepng} shows a Siemens PLC running in PLCSIM Advanced.

\begin{figure}[htbp]
    \centering
    \includegraphics[width=0.9\textwidth]{../attacks/down_up_program/image.png}
    \caption{Siemens PLC running in PLCSIM Advanced}
    \label{fig:cksdownupprogramimagepng}
\end{figure}

\textbf{Step 3} --- Connect the PLC to a WinCC HMI through an IE General connection:

Figure~\ref{fig:ksdownupprogramimage1png} shows the PLC connected to WinCC HMI via an IE General connection card.

\begin{figure}[htbp]
    \centering
    \includegraphics[width=0.9\textwidth]{../attacks/down_up_program/image-1.png}
    \caption{PLC connected to HMI WinCC via IE general connection card}
    \label{fig:ksdownupprogramimage1png}
\end{figure}

\textbf{Step 4} --- Create a PLC program, download it to the PLC, and run a test:

Figure~\ref{fig:ksdownupprogramimage3png} shows the PLC program in RUN mode with PLCSIM Advanced and WinCC HMI.

\begin{figure}[htbp]
    \centering
    \includegraphics[width=0.9\textwidth]{../attacks/down_up_program/image-3.png}
    \caption{PLC program in running mode with PLCSIM Advanced and HMI WinCC}
    \label{fig:ksdownupprogramimage3png}
\end{figure}

The screenshot shows the PLC program in RUN mode, the PLC instance in PLCSIM Advanced, and WinCC HMI displaying PLC variable values.

However, when observing packets in Wireshark, Wireshark could not parse S7comm payloads. Investigation showed that S7-1500 uses the S7comm-plus protocol and cannot be forced to use classic S7comm. According to \url{https://industrialmonitordirect.com/blogs/knowledgebase/s7comm-vs-s7comm-plus-switching-protocol-in-s7-1500-hmi-communication}, enabling \texttt{Permit access with PUT/GET communication from remote partners} only allows legacy clients to read selected data; Wireshark still cannot dissect S7comm-plus payloads.

Figure~\ref{fig:ksdownupprogramimage6png} shows Wireshark unable to parse S7comm-plus packets.

\begin{figure}[htbp]
    \centering
    \includegraphics[width=0.9\textwidth]{../attacks/down_up_program/image-6.png}
    \caption{Wireshark unable to parse S7comm plus packets}
    \label{fig:ksdownupprogramimage6png}
\end{figure}

Figure~\ref{fig:ksdownupprogramimage4png} shows Wireshark identifying only the outer COTP layer of S7comm-plus packets.

\begin{figure}[htbp]
    \centering
    \includegraphics[width=0.9\textwidth]{../attacks/down_up_program/image-4.png}
    \caption{Wireshark identifying only the outer COTP layer of S7comm plus packets}
    \label{fig:ksdownupprogramimage4png}
\end{figure}

Wireshark can identify the outer COTP envelope but cannot parse the S7comm-plus PDU carried inside it.

For Upload/Download rule tuning, this project therefore relies on documented sample captures from \url{https://wiki.wireshark.org/S7comm} rather than on lab-generated block-transfer PCAPs.

\end{document}
